\documentclass[10pt,aspectratio=169]{beamer}
\usepackage{bma01}

\title[Funciones por partes]{Funciones definidas por partes}
\subtitle{Valor absoluto, signo, escalón unitario y máximo entero\\[2pt]\small Semana 1 -- Sesión 2}

\begin{document}

\begin{frame}
  \titlepage
\end{frame}

\begin{frame}{Objetivos de la sesión}
  \begin{enumerate}
    \item Definir, evaluar y graficar funciones \resaltar{definidas por partes}.
    \item Manejar el \resaltar{valor absoluto}: propiedades, desigualdades y gráficas.
    \item Conocer las funciones \resaltar{signo}, \resaltar{escalón unitario} (Heaviside) y \resaltar{máximo entero}.
    \item Modelar situaciones con tarifas, interruptores y señales mediante estas funciones.
  \end{enumerate}
  \vspace{1mm}
  {\small\textbf{Competencia 1 del sílabo:} Construye modelos funcionales trabajando en equipo para resolver problemas de ciencia e ingeniería.}
  \begin{block}{Plan (2 horas; incluye verificación previa de 5 min)}
    \begin{tabular}{ll}
      0:00 -- 0:15 & Funciones por partes: definición y gráficas\\
      0:15 -- 0:50 & Valor absoluto: propiedades, ecuaciones e inecuaciones, gráficas\\
      0:50 -- 1:10 & Signo y escalón unitario\\
      1:10 -- 1:40 & Máximo entero: propiedades, gráficas, ecuaciones\\
      1:40 -- 2:00 & Actividad colaborativa, resumen y cierre metacognitivo
    \end{tabular}
  \end{block}
\end{frame}

\begin{frame}{Aula invertida: preparación y verificación}
  \small
  \begin{previa}[antes de la sesión]
    \begin{enumerate}
      \item Lea sobre valor absoluto, signo, escalón unitario y máximo entero.
      \item Grafique $|x|$, $\sgn x$ y $\ent{x}$, y reescriba $f(x)=|x-1|+|x+2|$ sin valor absoluto.
      \item Pregunta guía: ¿cuánto valen $\ent{-1.5}$ y $\ent{2}$?
    \end{enumerate}
  \end{previa}
  \begin{block}{Verificación al inicio (5 min, en pares)}
    \begin{enumerate}
      \item Calcule $\ent{-1.5}$ y $\ent{2}$.
      \item Resuelva $|x-3|<2$.
    \end{enumerate}
    Los pares discuten y votan; el docente aclara solo lo que genera dudas y ajusta el ritmo de la sesión.
  \end{block}
\end{frame}

\section{Funciones por partes}

\begin{frame}{Funciones definidas por partes}
  \begin{columns}
    \begin{column}{0.52\textwidth}
      Una función puede estar dada por \textbf{fórmulas distintas en distintos intervalos}:
      \[ f(x)=\begin{cases}x+2,&x<-1\\ x^2,&-1\le x\le1\\ 2-x,&x>1\end{cases} \]
      \begin{itemize}
        \item Para evaluar, primero se ubica el intervalo: $f(-3)=-1$, $f(\frac12)=\frac14$, $f(4)=-2$.
        \item En la gráfica: punto \textbf{lleno} $\bullet$ si el extremo pertenece; \textbf{vacío} $\circ$ si no.
      \end{itemize}
      \begin{observacion}
        Los intervalos no deben solaparse (si no, $f$ no estaría bien definida).
      \end{observacion}
    \end{column}
    \begin{column}{0.48\textwidth}
      \centering
      \begin{tikzpicture}
        \begin{axis}[bma, xmin=-4, xmax=4, ymin=-2.5, ymax=2.2, width=7cm, height=5.5cm, xtick={-1,1}, ytick={1}]
          \addplot[curva, domain=-4:-1] {x+2};
          \addplot[curva, domain=-1:1] {x^2};
          \addplot[curva, domain=1:4] {2-x};
          \addplot[only marks, mark=*, uniAzul] coordinates {(-1,1) (1,1)};
        \end{axis}
      \end{tikzpicture}

      \small Aquí las piezas ``empalman'': en $x=\pm1$ los valores coinciden.
    \end{column}
  \end{columns}
\end{frame}

\section{Valor absoluto}

\begin{frame}{La función valor absoluto}
  \begin{columns}
    \begin{column}{0.55\textwidth}
      \begin{definicion}
        \[ |x|=\begin{cases}x,&x\ge0\\-x,&x<0\end{cases}\qquad\text{equivalentemente } |x|=\sqrt{x^2}. \]
        $|x-a|$ es la \textbf{distancia} entre $x$ y $a$ en la recta real.
      \end{definicion}
      \begin{proposicion}
        \begin{enumerate}
          \item $|x|\ge0$; \ $|x|=0\iff x=0$; \ $|xy|=|x||y|$.
          \item $|x+y|\le|x|+|y|$ \ (desigualdad triangular).
          \item $|x|<a\iff-a<x<a$ \ ($a>0$).
          \item $|x|>a\iff x<-a$ ó $x>a$.
          \item $\big||x|-|y|\big|\le|x-y|$.
        \end{enumerate}
      \end{proposicion}
    \end{column}
    \begin{column}{0.45\textwidth}
      \centering
      \begin{tikzpicture}
        \begin{axis}[bma, xmin=-3, xmax=3, ymin=-0.5, ymax=3, width=6.5cm, height=4.5cm, xtick={-2,-1,1,2}, ytick={1,2}]
          \addplot[curva, domain=-3:3] {abs(x)};
        \end{axis}
      \end{tikzpicture}
      \vspace{2mm}

      \begin{tikzpicture}[scale=0.9]
        \draw[->] (-0.5,0) -- (6,0);
        \fill[uniAzul!30] (1.5,-0.12) rectangle (4.5,0.12);
        \draw[thick] (3,0.15) -- (3,-0.15) node[below] {\scriptsize $a$};
        \draw (1.5,0) node[below=3pt] {\scriptsize $a-\delta$};
        \draw (4.5,0) node[below=3pt] {\scriptsize $a+\delta$};
        \draw[uniAzul, thick] (1.5,0) circle (2pt); \draw[uniAzul, thick] (4.5,0) circle (2pt);
        \node[uniAzul] at (3,0.5) {\scriptsize $|x-a|<\delta$};
      \end{tikzpicture}

      \small Este tipo de intervalo será esencial para la definición de límite.
    \end{column}
  \end{columns}
\end{frame}

\begin{frame}{Ecuaciones e inecuaciones con valor absoluto}
  \begin{ejemplo}[1]
    $|2x-3|=5\iff2x-3=5$ ó $2x-3=-5\iff x=4$ ó $x=-1$.
  \end{ejemplo}
  \begin{ejemplo}[2]
    $|x-2|<0.1\iff1.9<x<2.1$ \quad (los $x$ a distancia menor que $0.1$ de $2$).
  \end{ejemplo}
  \begin{ejemplo}[3: análisis por casos]
    $|x-1|+|x+2|\le5$. Puntos críticos: $x=-2$, $x=1$.
    \begin{itemize}
      \item $x<-2$: \ $-(x-1)-(x+2)=-2x-1\le5\Rightarrow x\ge-3$. \ Solución $[-3,-2)$.
      \item $-2\le x<1$: \ $-(x-1)+(x+2)=3\le5$: \ siempre. \ Solución $[-2,1)$.
      \item $x\ge1$: \ $2x+1\le5\Rightarrow x\le2$. \ Solución $[1,2]$.
    \end{itemize}
    Conjunto solución: $[-3,2]$. \ (Interpretación: suma de distancias a $1$ y a $-2$ menor o igual que $5$.)
  \end{ejemplo}
\end{frame}

\begin{frame}{Gráficas con valor absoluto}
  \begin{columns}
    \begin{column}{0.33\textwidth}
      \centering
      \begin{tikzpicture}
        \begin{axis}[bma, xmin=-1.5, xmax=4.5, ymin=-2.3, ymax=3, width=5cm, height=4.5cm, xtick={1,3}, ytick={-1,1}, samples=100, domain=-1.2:4.2]
          \addplot[gray, dashed] {x^2-4*x+3};
          \addplot[curva2] {abs(x^2-4*x+3)};
        \end{axis}
      \end{tikzpicture}\\
      \small $y=|x^2-4x+3|$: se refleja la parte negativa.
    \end{column}
    \begin{column}{0.33\textwidth}
      \centering
      \begin{tikzpicture}
        \begin{axis}[bma, xmin=-4.5, xmax=4.5, ymin=-2.3, ymax=3, width=5cm, height=4.5cm, xtick={-3,-1,1,3}, ytick={-1,1}, samples=100, domain=-4.2:4.2]
          \addplot[gray, dashed, domain=-1.2:4.2] {x^2-4*x+3};
          \addplot[curva] {x^2-4*abs(x)+3};
        \end{axis}
      \end{tikzpicture}\\
      \small $y=f(|x|)$: se copia la parte $x\ge0$ simétricamente.
    \end{column}
    \begin{column}{0.33\textwidth}
      \centering
      \begin{tikzpicture}
        \begin{axis}[bma, xmin=-3.5, xmax=3.5, ymin=-0.5, ymax=5, width=5cm, height=4.5cm, xtick={-2,1}, ytick={3}, samples=100, domain=-3.5:3.5]
          \addplot[uniVerde, very thick] {abs(x-1)+abs(x+2)};
          \addplot[gray, dashed] {5};
        \end{axis}
      \end{tikzpicture}\\
      \small $y=|x-1|+|x+2|$: poligonal con ``codos'' en $-2$ y $1$.
    \end{column}
  \end{columns}
  \begin{observacion}
    $|f(x)|$ y $f(|x|)$ son transformaciones distintas. La gráfica de la derecha resuelve visualmente el Ejemplo 3: está por debajo
    de $y=5$ en $[-3,2]$.
  \end{observacion}
\end{frame}

\section{Signo y escalón unitario}

\begin{frame}{La función signo}
  \begin{columns}
    \begin{column}{0.55\textwidth}
      \begin{definicion}
        \[ \sgn(x)=\begin{cases}1,&x>0\\0,&x=0\\-1,&x<0\end{cases} \]
      \end{definicion}
      \begin{proposicion}
        \begin{itemize}
          \item $|x|=x\,\sgn(x)$; \ $x=|x|\sgn(x)$.
          \item $\sgn(xy)=\sgn(x)\sgn(y)$.
          \item $\sgn(x)=\dfrac{x}{|x|}$ para $x\ne0$.
        \end{itemize}
      \end{proposicion}
      \begin{ejemplo}
        $f(x)=\sgn(x^2-1)$ vale $1$ si $|x|>1$, $0$ si $x=\pm1$ y $-1$ si $|x|<1$.
      \end{ejemplo}
    \end{column}
    \begin{column}{0.45\textwidth}
      \centering
      \begin{tikzpicture}
        \begin{axis}[bma, xmin=-3, xmax=3, ymin=-1.6, ymax=1.6, width=6.3cm, height=4.6cm, xtick={-2,-1,1,2}, ytick={-1,1}]
          \addplot[curva] coordinates {(-3,-1) (0,-1)};
          \addplot[curva] coordinates {(0,1) (3,1)};
          \addplot[only marks, mark=o, uniAzul, fill=white] coordinates {(0,-1) (0,1)};
          \addplot[only marks, mark=*, uniAzul] coordinates {(0,0)};
        \end{axis}
      \end{tikzpicture}

      \small Salto de altura $2$ en el origen.
    \end{column}
  \end{columns}
\end{frame}

\begin{frame}{Escalón unitario (función de Heaviside)}
  \begin{columns}
    \begin{column}{0.55\textwidth}
      \begin{definicion}
        \[ u(t)=\begin{cases}0,&t<0\\1,&t\ge0\end{cases}\qquad u_a(t)=u(t-a). \]
      \end{definicion}
      Modela un \textbf{interruptor}: algo que se ``enciende'' en $t=a$.
      \begin{itemize}
        \item Pulso rectangular en $[a,b)$: \ $u(t-a)-u(t-b)$.
        \item ``Encender'' $g$ desde $t=a$: \ $g(t)\,u(t-a)$.
      \end{itemize}
      \begin{ejemplo}[circuito]
        Se aplica un voltaje de $12$ V entre $t=2$ s y $t=5$ s: \ $V(t)=12\big[u(t-2)-u(t-5)\big]$.
      \end{ejemplo}
    \end{column}
    \begin{column}{0.45\textwidth}
      \centering
      \begin{tikzpicture}
        \begin{axis}[bma, xlabel={$t$}, ylabel={$V$}, xmin=0, xmax=7, ymin=-1, ymax=14, width=6.3cm, height=4.3cm, xtick={2,5}, ytick={12}]
          \addplot[curva] coordinates {(0,0) (2,0)};
          \addplot[curva] coordinates {(2,12) (5,12)};
          \addplot[curva] coordinates {(5,0) (7,0)};
          \addplot[only marks, mark=*, uniAzul] coordinates {(2,12) (5,0)};
          \addplot[only marks, mark=o, uniAzul, fill=white] coordinates {(2,0) (5,12)};
        \end{axis}
      \end{tikzpicture}

      \begin{tikzpicture}
        \begin{axis}[bma, xlabel={$t$}, xmin=-1, xmax=6, ymin=-1.3, ymax=1.3, width=6.3cm, height=3.6cm, xtick={1}, ytick={1}, samples=150]
          \addplot[curva2, domain=-1:1] {0};
          \addplot[curva2, domain=1:6] {sin(deg(2*(x-1)))};
        \end{axis}
      \end{tikzpicture}

      \small $\sen\big(2(t-1)\big)\,u(t-1)$: señal que comienza en $t=1$.
    \end{column}
  \end{columns}
\end{frame}

\section{Máximo entero}

\begin{frame}{La función máximo entero}
  \begin{columns}
    \begin{column}{0.55\textwidth}
      \begin{definicion}
        $\ent{x}$ es el mayor entero menor o igual que $x$:
        \[ \ent{x}=n\iff n\le x<n+1,\quad n\in\Z. \]
      \end{definicion}
      \begin{ejemplo}
        $\ent{2.7}=2$, \ $\ent{5}=5$, \ $\ent{-0.3}=-1$, \ $\ent{-2.5}=-3$, \ $\ent{\pi}=3$.
      \end{ejemplo}
      \begin{proposicion}
        \begin{enumerate}
          \item $\ent{x}\le x<\ent{x}+1$, \ es decir \ $x-1<\ent{x}\le x$.
          \item $\ent{x+n}=\ent{x}+n$ si $n\in\Z$.
          \item $\ent{x}=x\iff x\in\Z$.
          \item $\ent{x}\ge n\iff x\ge n$ \ ($n\in\Z$).
        \end{enumerate}
      \end{proposicion}
    \end{column}
    \begin{column}{0.45\textwidth}
      \centering
      \begin{tikzpicture}
        \begin{axis}[bma, xmin=-3.3, xmax=3.6, ymin=-3.5, ymax=3.5, width=6.5cm, height=6.3cm, xtick={-3,...,3}, ytick={-3,...,3}]
          \foreach \n in {-3,...,2} {
            \addplot[curva] coordinates {(\n,\n) (\n+1,\n)};
          }
          \addplot[only marks, mark=*, uniAzul] coordinates {(-3,-3) (-2,-2) (-1,-1) (0,0) (1,1) (2,2)};
          \addplot[only marks, mark=o, uniAzul, fill=white] coordinates {(-2,-3) (-1,-2) (0,-1) (1,0) (2,1) (3,2)};
          \addplot[gray, dashed, domain=-3:3.5] {x};
        \end{axis}
      \end{tikzpicture}

      \small Escalera: ``saltos'' de $1$ en cada entero. Siempre por debajo (o sobre) $y=x$.
    \end{column}
  \end{columns}
\end{frame}

\begin{frame}{Más gráficas con máximo entero}
  \begin{columns}
    \begin{column}{0.5\textwidth}
      \centering
      \begin{tikzpicture}
        \begin{axis}[bma, xmin=-1.3, xmax=3.3, ymin=-0.4, ymax=1.4, width=6.5cm, height=4.3cm, xtick={-1,...,3}, ytick={1}]
          \foreach \n in {-1,...,2} {
            \addplot[uniVerde, very thick] coordinates {(\n,0) (\n+1,1)};
          }
          \addplot[only marks, mark=*, uniVerde] coordinates {(-1,0) (0,0) (1,0) (2,0)};
          \addplot[only marks, mark=o, uniVerde, fill=white] coordinates {(0,1) (1,1) (2,1) (3,1)};
        \end{axis}
      \end{tikzpicture}\\
      \small $y=x-\ent{x}$ \ (parte fraccionaria; ``diente de sierra''). Rango $[0,1)$; periódica de periodo $1$.
    \end{column}
    \begin{column}{0.5\textwidth}
      \centering
      \begin{tikzpicture}
        \begin{axis}[bma, xmin=-0.3, xmax=2.3, ymin=-0.5, ymax=4.6, width=6.5cm, height=4.3cm, xtick={0.5,1,1.5,2}, ytick={1,2,3,4}]
          \foreach \n in {0,...,3} {
            \addplot[curva2] coordinates {(\n/2,\n) (\n/2+0.5,\n)};
          }
          \addplot[only marks, mark=*, uniGranate] coordinates {(0,0) (0.5,1) (1,2) (1.5,3)};
          \addplot[only marks, mark=o, uniGranate, fill=white] coordinates {(0.5,0) (1,1) (1.5,2) (2,3)};
        \end{axis}
      \end{tikzpicture}\\
      \small $y=\ent{2x}$: escalones de ancho $\frac12$.
    \end{column}
  \end{columns}
  \begin{ejemplo}[ecuación]
    $\ent{2x-1}=3\iff3\le2x-1<4\iff2\le x<\frac52$.
  \end{ejemplo}
\end{frame}

\begin{frame}{Aplicación: tarifas escalonadas}
  \begin{columns}
    \begin{column}{0.55\textwidth}
      \begin{ejemplo}[estacionamiento]
        Un estacionamiento cobra S/ 3 por la primera hora o fracción y S/ 2 por cada hora adicional o fracción. Si $t>0$ son las horas:
        \[ C(t)=\begin{cases}3,&0<t\le1\\ 3+2\left\lceil t-1\right\rceil,&t>1\end{cases} \]
        donde $\lceil x\rceil=-\ent{-x}$ (redondeo hacia arriba).
        \medskip

        $C(2.25)=3+2\cdot2=7$ soles.
      \end{ejemplo}
      \begin{observacion}
        Las funciones escalonadas modelan magnitudes \textbf{cuantizadas}: tarifas, niveles de energía de un átomo, bits.
      \end{observacion}
    \end{column}
    \begin{column}{0.45\textwidth}
      \centering
      \begin{tikzpicture}
        \begin{axis}[bma, xlabel={$t$}, ylabel={$C$}, xmin=0, xmax=4.4, ymin=0, ymax=11, width=6.3cm, height=5.3cm, xtick={1,2,3,4}, ytick={3,5,7,9}]
          \addplot[curva] coordinates {(0,3) (1,3)};
          \addplot[curva] coordinates {(1,5) (2,5)};
          \addplot[curva] coordinates {(2,7) (3,7)};
          \addplot[curva] coordinates {(3,9) (4,9)};
          \addplot[only marks, mark=*, uniAzul] coordinates {(1,3) (2,5) (3,7) (4,9)};
          \addplot[only marks, mark=o, uniAzul, fill=white] coordinates {(0,3) (1,5) (2,7) (3,9)};
        \end{axis}
      \end{tikzpicture}
    \end{column}
  \end{columns}
\end{frame}

\begin{frame}{Actividad colaborativa: Tarifa de un estacionamiento}
  \small
  \begin{activa}[Estudio de casos / ABP]
    Un estacionamiento cobra S/ 4 la primera hora o fracción, S/ 3 por cada hora o fracción adicional, con un máximo de S/ 20 por día. \textbf{(a)} Modele el costo $C(t)$ con $\ent{\ }$ y por partes. \textbf{(b)} Grafíquelo señalando puntos abiertos y cerrados. \textbf{(c)} Proponga una tarifa alternativa y defienda por qué es más justa.
  \end{activa}
  \begin{itemize}
    \item \textbf{Organización:} equipos de 3 a 4 con roles rotativos (coordinador, redactor, verificador, portavoz); 12--15 minutos.
    \item \textbf{Producto:} Modelo, gráfica y argumento de la propuesta tarifaria (media página).
    \item \textbf{Retroalimentación:} el docente recorre los equipos, hace preguntas y no da la solución hasta la puesta en común.
  \end{itemize}
\end{frame}

\section{Ejercicios}

\begin{frame}{Ejercicios propuestos}
  \small
  \begin{ejercicio}
    \begin{enumerate}
      \item Grafique $f(x)=\begin{cases}\sqrt{-x},&x<0\\ 1-x^2,&0\le x<2\\ x-5,&x\ge2\end{cases}$ \ y halle su rango.
      \item Resuelva: (a) $|3x+2|=|x-4|$, \ (b) $|x^2-4|<3$, \ (c) $|x|+|x-3|>5$, \ (d) $\left|\dfrac{x-1}{x+1}\right|\le1$.
      \item Grafique: (a) $y=\big||x|-2\big|$, \ (b) $y=x|x|$, \ (c) $y=\sgn(\sen x)$, \ (d) $y=x\,\sgn(x-1)$.
      \item Exprese $f(t)=\begin{cases}0,&t<1\\ t,&1\le t<3\\ 0,&t\ge3\end{cases}$ \ usando escalones unitarios.
      \item Grafique: (a) $y=\ent{x}+\ent{-x}$, \ (b) $y=\ent{x^2}$ en $[0,2]$, \ (c) $y=x\ent{x}$ en $[-2,2]$, \ (d) $y=\ent{\sen x}$.
      \item Resuelva: (a) $\ent{\frac{x+1}{2}}=-1$, \ (b) $\ent{x}^2-\ent{x}-2=0$, \ (c) $\ent{|x|}<2$.
      \item Pruebe que $\ent{x}+\ent{x+\frac12}=\ent{2x}$ (identidad de Hermite para $n=2$).
    \end{enumerate}
  \end{ejercicio}
\end{frame}

\begin{frame}{Metacognición y uso responsable de la IA}
  \small
  \begin{columns}[t]
    \begin{column}{0.5\textwidth}
      \begin{metacog}[cierre, 3 min]
        \begin{enumerate}
          \metacomun
          \item ¿En qué punto me equivoqué al tratar los extremos abiertos y cerrados de cada tramo?
        \end{enumerate}
      \end{metacog}
    \end{column}
    \begin{column}{0.5\textwidth}
      \begin{iabox}
        Pida a una IA graficar $\ent{x}+\ent{-x}$. Verifique a mano en $x=0.5,\,1,\,-1.5$ si la gráfica respeta los puntos abiertos y cerrados.\\[2pt]
        \textit{Regla:} primero intente usted; luego use la herramienta para contrastar, nunca para sustituir su razonamiento. Anote el cambio en su bitácora.
      \end{iabox}
    \end{column}
  \end{columns}
\end{frame}

\begin{frame}{Resumen}
  \begin{itemize}
    \item Por partes: identificar el intervalo; extremos con $\bullet$ (incluido) o $\circ$ (excluido).
    \item $|x-a|$ = distancia; $|x-a|<\delta\iff a-\delta<x<a+\delta$; desigualdad triangular.
    \item $|f(x)|$ refleja la parte negativa; $f(|x|)$ simetriza la parte derecha.
    \item $\sgn(x)$ y $u(t)$: saltos; $u(t-a)-u(t-b)$ es un pulso.
    \item $\ent{x}=n\iff n\le x<n+1$; escalera; $x-\ent{x}\in[0,1)$.
  \end{itemize}
  \vspace{3mm}
  \textbf{Próxima sesión:} funciones monótonas, periódicas, inyectivas, sobreyectivas y biyectivas.
\end{frame}


\begin{frame}{Evaluación y trabajo posterior}
  \small
  \begin{evalbox}[Práctica del miércoles de la semana 1: Prueba de entrada]
    \textbf{Cobertura:} conocimientos previos de precálculo. \textbf{Naturaleza:} diagnóstica, individual, sin nota para el promedio; sirve para ajustar la nivelación y los grupos de trabajo.
  \end{evalbox}
  \begin{itemize}
    \item \criterios
    \item \declaracionia
    \item \textbf{Portafolio:} guarde el producto de la actividad, la bitácora de metacognición y los ejercicios propuestos de esta sesión; son evidencia de progreso y base de la retroalimentación.
  \end{itemize}
\end{frame}

\end{document}
